

% ------------------------------------------------------------
\begin{frame}{Related Literature}
\hypertarget{src:literature}{}
\footnotesize
Four literatures each name a piece of the phenomenon; none estimates what litigation does to a race.
\begin{itemize}\itemsep1pt
  \item \emph{Judicialization of competition.} A comparative literature documents how courts come to settle the terms of political competition \ns{\citep{tate1995,hirschl2008,sieder2020}}, and in Brazil how the TSE rewrites electoral rules through adjudication \ns{\citep{marchetticortez2009}}. This work is descriptive; we estimate the effect on the contest.

  \item \emph{Electoral governance.} A second literature defines electoral governance as rule-making, rule application and adjudication \ns{\citep{mozaffar2002}} and, for Brazil, studies the fusion of all three in one branch as a problem of institutional design \ns{\citep{marchetti2008,marchetti2012}}. We instead treat litigation intensity as a shock that candidates and voters respond to.

  \item \emph{Lawsuits as strategy versus accountability.} A third body of work reads litigation either as a strategic weapon aimed at rivals \ns{\citep{chin2025,nakaguma2025}} or as a channel through which judicial decisions convey information \ns{\citep{assumpcao2019}}, with campaign money as a driver \ns{\citep{mancuso2023}}. These designs are candidate-level: who files, who is targeted, and how a conviction moves the convicted candidate's vote share. The one test on polls rules out drops in the defendant's vote intentions larger than two to five points \ns{\citep{chin2025}}. We estimate the effect on the race itself.

  \item \emph{Shift-share identification.} A methodological literature sets out when Bartik instruments identify, through exogenous shares \ns{\citep{goldsmith2020}} or many exogenous shocks \ns{\citep{borusyak2022}}, and how to do inference when units with similar shares have correlated residuals \ns{\citep{adao2019}}. Applications run from import exposure \ns{\citep{autor2013}} to leave-one-out legal-topic flows \ns{\citep{ash2025}}. The shock route needs many diffuse shocks; litigation data have a few concentrated ones, so we use the share-exogeneity (GPS) route.
\end{itemize}
\vfill
{\footnotesize\centering The first three literatures characterize the phenomenon; the fourth lets us estimate its effect.\par}
\end{frame}
